\section{Validating Global Performance Objectives}\label{sec:verification}
As we saw in the previous section, guardrails are local mechanisms that monitor a run-time condition and trigger a corrective action when that condition fails. But the property that the system designer ultimately cares about is usually more global. For example, in CBMM, the guardrail checks whether the ordering induced by observed accesses agrees with the ordering induced by the stored benefit estimates. That is a natural run-time check, but it is not itself the real objective. The real performance objective is agreement with an oracle promotion rule defined using the true deployment-time cost and benefit model, such as:
\[
\forall P.\;
\mathsf{Promote}_{\pi}(P)
\iff
\bigl(b^\star(P) > c^\star(P)\bigr),
\]
where $b^\star(P)$ and $c^\star(P)$ denote the true deployment-time benefit and cost of promoting page region $P$, respectively. 

More generally, a guardrail condition is useful only if maintaining or restoring it is enough to support the higher-level performance goal that motivated it.  To reason about this connection formally, we introduce a lightweight verification interface. The key idea is to separate \emph{what a guardrail monitors and repairs} from \emph{what the system ultimately wants to guarantee}. The former is given by the guardrail itself, as defined in Section~\ref{sec:framwework}; the latter is expressed as a global performance objective, which may depend on quantities that are not directly observable at run time. Verification relates the two under an abstract model of policy--system interaction.

\paragraph{Verification Interface}

Our verification interface required three pieces from the user. The first is an abstract model
$ 
M(\sysstate,\ppol{\policy}{\args},\sysstate')
$ 
that over-approximates the relevant policy-system dynamics: if the concrete system can step from $\sysstate$ to $\sysstate'$ under policy $\ppol{\policy}{\args}$ (written $\TransitionsVia{\sysstate}{\policy}{\sysstate}$), then so can $M$, although $M$ may permit additional behaviors to simplify reasoning. The second is a \emph{local guarantee}
$ 
\varphi(\sysstate,\ppol{\policy}{\args}),
$
which captures the property that a particular guardrail is intended to maintain or restore.\footnote{This per-guardrail property is annotated via an \texttt{ensures} clause in \lang.} The third is a \emph{global performance objective}
 $
\Phi(\sysstate,\ppol{\policy}{\args}),
$ (\perfsafety in Figure~\ref{fig:framework-overview})
which formalizes the higher-level property the system designer ultimately cares about.
This \perfsafety property, $\Phi$, can be over the long-horizon impact of a policy, or the net behavior of a subsystem, or some other system-wide objective\footnote{This global property is annotated via an \texttt{assert} clause in \lang.}.
\paragraph{Methodology.}
The verification task has three parts. First, for each guardrail, we check that its monitored condition and corrective action establish the appropriate \emph{local guarantee} defined by $\varphi$. Second, we check that the local guarantees of the installed guardrails are collectively strong enough to imply the desired \emph{global performance objective} $\Phi$. Third, when multiple guardrails may be enabled on the same policy, we must check that the action of one guardrail does not invalidate the guarantees maintained by the others. 

\paragraph{Local proof obligations.}
For each guardrail $G_i = (\property{i},\action{i},\enabled_i)$ and chosen local guarantee $\varphi_i$, the verifier checks the following two obligations:
\[ \small
\begin{array}{l@{\hspace{0.1em}}l@{\;}l}
\multicolumn{3}{l}{
M(s, \pi, s') \Rightarrow
\hspace{4em}\textit{(i.e. under the system model)}
} \\[0.7em]
& E_i(\pi) \wedge 
\phantom{\neg}
P_i(s')
\Rightarrow 
\varphi_i(\pi',s')
\wedge
(\pi' = \pi)
&
\textsc{Maintain}_{i}(\pi, s', \pi')
\\[0.8em]
& E_i(\pi) \wedge \neg P_i(s')
\Rightarrow
\varphi_i(\pi',s')
\wedge
(\pi' = A_i(s',\pi))
&
\textsc{Recover}_i(\pi, s', \pi')
% \\[0.8em]
% E_i(\pi) \wedge \varphi_i(\pi,s)
% \Rightarrow
% \varphi_i(A_i(s,\pi),s)
% &
% \textsc{Stable}_i
\end{array}
\]



The \textsc{Maintain} obligation says that if guardrail $G_i$ is enabled and its monitored condition already holds, then one model step lands in a state satisfying the local guarantee $\varphi_i$. The \textsc{Recover} obligation covers the complementary case: if the monitored condition fails, then applying the corrective action restores $\varphi_i$.
%Finally, \textsc{Stable} requires that the corrective action be non-destructive with respect to its own guarantee: if $\varphi_i$ already holds, then applying $A_i$ should preserve it.

\paragraph{Verifying Performance}
To verify that a guardrail $G_i$ ensures our \perfsafety condition $\Phi$, we must prove:
\[
% M(s, \pi, s') 
% \wedge
\textsc{Maintain}_i(\pi, s', \pi')
\wedge 
\textsc{Recover}_i(\pi, s', \pi')
\Rightarrow \Phi(s', \pi')
\]
The above obligation says that when a guardrail ensures their local proof obligations, the \perfsafety property. We further generalize the verification condition to systems guarded by sets of guardrails. Sound composition requires some additional \emph{concurrent noninterference} obligation~\cite{ewd117} which ensure that no guardrail ``interferes'' with the state of another. We formally proved this condition ensures soundness in the Lean theorem prover. The proof is reported in Appendix~\ref{appdx:verification-soundness}.
